\chapter{Astropy による時刻・単位・\mbox{天体データ処理}}
\label{chap:astropy}

Astropy\index{astropy@Astropy|textbf} は、単位付き物理量、時刻系、天球座標系、FITS、天文向け表を扱うライブラリである。単位、時刻系、座標系を数値とは別の文字列として管理すると、演算時に対応関係を失う可能性がある。本章では、これらの情報を数値と同じオブジェクトに保持し、変換、演算、表および FITS への保存を行う方法を扱う。

\section{数値と単位・時刻・座標の一元的な管理}

\subsection{数値への単位・時刻・座標情報の付与}

Astropy の中心的な発想は、「数値そのもの」ではなく「何の単位で、どの時刻系で、どの座標系なのか」までオブジェクトへ持たせることである。例えば 100 という値だけでは、100 ns なのか 100 ms なのか分からない。Astropy の \code{Quantity}\index{quantity@\code{Quantity}|textbf} を使えば、この曖昧さを減らせる。

\begin{lstlisting}[language=Python]
import astropy.units as u

time_window = 250 * u.ns
speed = 299792458 * u.m / u.s
distance = speed * time_window

print(distance)
print(distance.to(u.m))
\end{lstlisting}

この計算では、単位付きのまま距離が評価され、次のように表示される。

\begin{lstlisting}[numbers=none, xleftmargin=2\zw]
74.9481145 m
74.9481145 m
\end{lstlisting}

この例では、\code{250 * u.ns} という時点で「250 ns」という意味が値に埋め込まれている。式の途中で単位を落とさない限り、単位換算はライブラリが引き受ける。

数式で書けば、ここで行っているのは
\[
d = v \Delta t
\]
である。特に \(v = c\) とみなすなら、250 ns の時間窓に対応する距離 \(d = c \Delta t\) をそのまま安全に計算できる。

\subsection{時刻系・座標系を属性として保持する理由}

観測データでは、UTC、TAI、TT、MJD、JD、Unix time が混在することがある。文字列の分解だけで相互変換すると、時刻系のオフセットとうるう秒が計算に反映されず、時刻差を誤る可能性がある。Astropy の \code{Time} は、時刻形式と時刻系を属性として保持する。

\section{単位・時刻・座標・表・FITS の操作}

\subsection{\code{units} と \code{Quantity}}

単位付き計算の基本は \code{Quantity} である。

\begin{lstlisting}[language=Python]
import astropy.units as u

length = 3.2 * u.m
time = 12.5 * u.ns
velocity = length / time

print(velocity)
print(velocity.to(u.m / u.s))
\end{lstlisting}

単位変換の結果は次の通りである。

\begin{lstlisting}[numbers=none, xleftmargin=2\zw]
0.256 m / ns
256000000.0 m / s
\end{lstlisting}

\code{to()} は、指定した単位へ換算した \code{Quantity} を返す。\code{.value} は単位情報を除いた数値を返すため、単位を受け取らない関数へ渡す直前など、数値だけが必要な箇所に使用を限定する。

この計算は
\[
v = \frac{L}{t}
\]
に対応しており、具体的には
\[
v = \frac{3.2\ \mathrm{m}}{12.5\ \mathrm{ns}}
= 0.256\ \mathrm{m/ns}
= 2.56 \times 10^8\ \mathrm{m/s}
\]
である。コードの出力をこの換算式と照合すれば、桁と SI 接頭語の対応を検証できる。

\subsection{物理定数}

Astropy には物理定数も入っている。SciPy の \code{constants} と似ているが、こちらは単位付きで返る点が大きい。

\begin{lstlisting}[language=Python]
from astropy.constants import c, k_B

print(c)
print(k_B)
\end{lstlisting}

Astropy の定数オブジェクトを表示すると、値だけでなく単位と参照元も確認できる。

\begin{lstlisting}[numbers=none, xleftmargin=2\zw]
  Name   = Speed of light in vacuum
  Value  = 299792458.0
  Uncertainty  = 0.0
  Unit  = m / s
  Reference = CODATA 2022
  Name   = Boltzmann constant
  Value  = 1.380649e-23
  Uncertainty  = 0.0
  Unit  = J / K
  Reference = CODATA 2022
\end{lstlisting}

この表示から、定数の値、単位、不確かさ、参照する定数体系を確認できる。

\subsection{\code{Time} による時刻表現}

Astropy の \code{Time}\index{time@\code{Time}|textbf} は、複数のフォーマットと複数の時刻系をまたいで扱える。

\begin{lstlisting}[language=Python]
from astropy.time import Time

t = Time("2026-04-12 12:34:56.123456789", format="iso", scale="utc")

print(t.mjd)
print(t.unix)
print(t.tai.iso)
\end{lstlisting}

同じ時刻を複数の表現へ変換すると、次の値が得られる。

\begin{lstlisting}[numbers=none, xleftmargin=2\zw]
61142.52426068816
1775997296.1234567
2026-04-12 12:35:33.123
\end{lstlisting}

\code{format="iso"} は入力形式、\code{scale="utc"} は時刻系を表す。\code{mjd} と \code{unix} は同じ時刻の別形式を返し、\code{tai.iso} は時刻系を TAI へ変換して ISO 形式で返す。

たとえば、修正ユリウス日 MJD は
\[
\mathrm{MJD} = \mathrm{JD} - 2400000.5
\]
で定義される。Astropy の \code{Time} は、時刻形式と時刻系を属性として保持し、それらの変換を同じオブジェクト上で行う。

\subsection{時刻差と \code{TimeDelta}}

差分計算も \code{Time} のまま行える。\code{TimeDelta}\index{timedelta@\code{TimeDelta}|textbf} によって、差分の単位変換も一貫して扱える。

\begin{lstlisting}[language=Python]
from astropy.time import Time

t1 = Time("2026-04-12T12:00:00.000000100", format="isot", scale="utc")
t2 = Time("2026-04-12T12:00:00.000000320", format="isot", scale="utc")
dt = t2 - t1

print(f"{dt.to_value('s'):.9e}")
print(f"{dt.to_value('ns'):.3f}")
\end{lstlisting}

時刻差を秒と ns へ変換した結果は次の通りである。

\begin{lstlisting}[numbers=none, xleftmargin=2\zw]
2.200000182e-07
220.000
\end{lstlisting}

\code{TimeDelta} の既定の内部形式や表示は、入力形式と同じ単位とは限らない。そのため、出力では \code{to\_value("s")} や \code{to\_value("ns")} によって単位を明示する。秒境界や日付境界をまたぐ場合も、文字列を分解せずに差を計算できる。

ここで計算している時刻差は
\[
\Delta t = t_2 - t_1 = 220\ \mathrm{ns}
\]
である。\code{Time} どうしの減算は \code{TimeDelta} を返すため、時刻系に基づく差分と単位情報が結果に保持される。

\subsection{\code{SkyCoord} による座標変換}

Astropy の \code{coordinates} は、天球上の座標系を変換する。\code{SkyCoord}\index{skycoord@\code{SkyCoord}|textbf} はその中心となるクラスである。

\begin{lstlisting}[language=Python]
from astropy.coordinates import SkyCoord
import astropy.units as u

coord = SkyCoord(ra=150.0 * u.deg, dec=20.0 * u.deg, frame="icrs")
gal = coord.galactic

print(f"{gal.l.deg:.6f}")
print(f"{gal.b.deg:.6f}")
\end{lstlisting}

座標変換後の銀河経度と銀河緯度は次の通りである。

\begin{lstlisting}[numbers=none, xleftmargin=2\zw]
213.800977
50.263718
\end{lstlisting}

\code{ra} と \code{dec} には単位を付ける。これにより、入力角度が度またはラジアンのどちらで与えられたかが \code{SkyCoord} に保持される。

この操作は、赤道座標 \((\alpha, \delta)\) から銀河座標 \((l, b)\) への変換に対応する。天球上の回転そのものは自分で式を書かなくても、入力座標系と出力座標系を明示すれば Astropy が一貫して扱える。

\subsection{\code{Table} と \code{QTable}}

Astropy には表処理のために \code{Table}\index{table@\code{Table}|textbf} と \code{QTable}\index{qtable@\code{QTable}|textbf} がある。\code{QTable} は列へ単位を保持できる。

\begin{lstlisting}[language=Python]
from astropy.table import QTable
import astropy.units as u

table = QTable()
table["time"] = [0.0, 1.0, 2.0] * u.s
table["signal"] = [12.4, 10.1, 14.8] * u.mV
table["detector"] = ["A", "A", "B"]

print(table)
\end{lstlisting}

\code{QTable} を表示すると、列名だけでなく単位も一緒に確認できる。

\begin{lstlisting}[numbers=none, xleftmargin=2\zw]
time signal detector
 s    mV
---- ------ --------
0.0   12.4        A
1.0   10.1        A
2.0   14.8        B
\end{lstlisting}

\code{QTable} は \code{Quantity} 列を保持するため、単位付きの列をそのまま演算できる。一方、複雑な groupby、join、列演算を中心とする処理には pandas や Polars の機能を利用できる。単位とメタデータの保持、および必要な表操作の双方から中心となる表形式を選ぶ。

\subsection{Astropy Table と pandas の行き来}

Astropy の表と pandas の DataFrame は相互変換できるが、単位や時刻の情報が完全には残らないことがある。

\begin{lstlisting}[language=Python]
df = table.to_pandas()
print(df)
\end{lstlisting}

\code{to\_pandas()} で得た通常の pandas 列には、\code{Quantity} の単位属性はそのまま保持されない。変換前に列単位を記録し、変換後の数値へどの単位を対応させるかを明示する。マスク付き列と時刻列についても、欠損表現と型が変換前後で一致するかを確認する。

\subsection{\code{astropy.io.fits} による FITS の読み書き}

天文データでは FITS\index{fits@FITS|textbf} が標準的である。Astropy なら NumPy 配列やヘッダをそのまま扱える。

\begin{lstlisting}[language=Python]
import numpy as np
from astropy.io import fits

image = np.arange(9).reshape(3, 3)
hdu = fits.PrimaryHDU(image)
hdu.header["OBJECT"] = "test"
hdu.writeto("image.fits", overwrite=True)

with fits.open("image.fits") as hdul:
    data = hdul[0].data
    header = hdul[0].header

print(data.shape)
print(header["OBJECT"])
\end{lstlisting}

FITS を書いてから読み直すと、配列形状とヘッダ値は次のように確認できる。

\begin{lstlisting}[numbers=none, xleftmargin=2\zw]
(3, 3)
test
\end{lstlisting}

\code{fits.open()} は HDU のリストを返す。画像は \code{PrimaryHDU}、表は \code{BinTableHDU} などに格納される。\code{hdul.info()} で各 HDU の型、名前、次元を確認してから、対象のデータを読み込む。

\subsection{ASCII・ECSV 形式の読み込み}

Astropy の表入出力は FITS だけではない。\code{Table.read()} は ASCII、CSV、ECSV などにも対応している。

\begin{lstlisting}[language=Python]
from astropy.table import Table

tbl = Table.read("events.ecsv", format="ascii.ecsv")
print(tbl.colnames)
print(len(tbl))
\end{lstlisting}

\code{Table.read()} の第 1 引数は入力パス、\code{format="ascii.ecsv"} は ECSV として解釈する指定である。3 行で \code{time} と \code{signal} の 2 列を持つ表なら、\code{['time', 'signal']} と \code{3} が表示される。ECSV は、列の単位、データ型、表のメタデータをテキストとともに保存できるため、中間生成物の受け渡しに利用できる。保存時は \code{table.write("events.ecsv", format="ascii.ecsv", overwrite=True)} と書く。\code{overwrite=True} は同名ファイルの置換を許可する指定である。

\section{単位情報・時刻系・座標系を失わないための注意}

\subsection{\code{.value} を早く使い過ぎない}

\code{Quantity.value} は、単位を除いた NumPy 配列またはスカラーを返す。以後の計算では単位整合性が検査されないため、外部ライブラリへ渡す直前まで \code{Quantity} を保持し、\code{.to\_value("ns")} のように変換先の単位を明示する。

\subsection{時刻系を省略しない}

\code{Time("2026-04-12 12:00:00")} のように \code{format} と \code{scale} を省略すると、Astropy の推定と既定値に依存し、作成者が意図した入力形式と時刻系がコード上に残らない。観測データでは \code{format=} と \code{scale=} を明示し、UTC、TAI、TT などを区別する。

\subsection{座標変換に必要な観測条件}

赤道座標から銀河座標への変換だけなら比較的単純だが、地平座標へ変換するには観測地点や観測時刻が必要になる。座標系だけを変えているつもりでも、背後では時刻と場所が意味を持つことがある。

\subsection{DataFrame への変換に伴う情報の欠落}

Astropy Table を pandas や Polars へ変換すると、単位、マスク、時刻系、メタデータの一部は対応する表現を持たず、失われたり別の型へ変換されたりする。変換前後の列型とメタデータを比較し、必要な情報は別の列または付随するメタデータとして保存する。集計のために変換する場合も、計算後に単位を復元できる情報を残す。

\section{SciPy・pandas とのデータ受け渡し}

Astropy を使うべき場面は、数値とともに単位、時刻系、座標系、天文形式の情報を保持する必要がある場合である。代表的な用途を次に示す。

\begin{itemize}
\item 観測時刻を UTC、MJD、JD、Unix time の間で変換したいとき
\item ns、\(\mu\)s、mV、MeV などの単位を混ぜて計算したいとき
\item 天球座標を別の座標系へ変換したいとき
\item FITS 画像や FITS テーブルを直接読み書きしたいとき
\item 表データに単位やメタデータを持たせたまま運びたいとき
\end{itemize}

大規模表の集計や複雑な join では、pandas や Polars が必要な API と性能を備えている場合がある。その場合は、Astropy で単位と時刻を解釈し、変換時に失われる情報を記録してから必要な列だけを表処理ライブラリへ渡す。集計結果を Matplotlib や SciPy へ渡すときも、数値に対応する単位と時刻系を併記する。

全体の導入は Getting Started~\cite{astropy_getting_started} と User Guide~\cite{astropy_user_guide}、単位は Units and Quantities~\cite{astropy_units}、時刻は Time and Dates~\cite{astropy_time}、座標は Coordinates~\cite{astropy_coordinates}、表は Tables~\cite{astropy_table}、FITS は FITS File Handling~\cite{astropy_fits} に記載されている。本章は、そのうち観測データ解析で用いる機能を扱った。

\section{単位付き観測データの処理}

\begin{enumerate}
\item \textbf{単位を保持した時間窓の比較}。
\code{window = 250 * u.ns} とし、\code{dt = np.array([120, 220, 250, 251, 800]) * u.ns} の各要素が時間窓以内かを判定せよ。真になる要素は 120、220、250 ns の 3 個である。境界の 250 ns を含めるために \code{<} ではなく \code{<=} を使う理由と、\code{u.ns} を付けずに比較した場合の問題を説明せよ。

\item \textbf{秒境界をまたぐ時刻差の計算}。
\code{2026-04-12T11:59:59.999999900} と \code{2026-04-12T12:00:00.000000120} を UTC の \code{Time} として作り、差を ns で表示せよ。結果は約 220 ns になり、250 ns 以内である。\code{format="isot"}、\code{scale="utc"}、\code{to\_value("ns")} の役割を説明せよ。

\item \textbf{UTC と TAI による同一時刻の表現}。
同じ \code{Time} オブジェクトについて \code{utc.iso} と \code{tai.iso} を表示し、文字列表現が異なっても同じ瞬間を表していることを確認せよ。文字列を直接引き算する方法が不適切である理由を、時刻系とうるう秒の観点から説明せよ。

\item \textbf{単位付き表の ECSV と FITS への保存}。
\code{time}、\code{signal}、\code{detector} を持つ \code{QTable} を作り、\code{time} に ns、\code{signal} に mV を付けよ。ECSV へ保存して読み戻し、各列の単位と行数を表示せよ。さらに同じ信号値を FITS 画像として保存し、画像 HDU のヘッダへ \code{BUNIT="mV"} を記録する。\code{fits.open()} で形状、\code{BUNIT}、信号値を読み戻し、ECSV の列単位と一致することを確認せよ。\code{overwrite=True} を付ける場合の挙動も説明せよ。

\item \textbf{座標系を明示した変換}。
赤経 150 度、赤緯 20 度を ICRS の \code{SkyCoord} として作り、銀河座標へ変換せよ。銀河経度は約 213.801 度、銀河緯度は約 50.264 度になる。\code{ra}、\code{dec}、\code{frame} の意味を述べ、地平座標への変換では観測時刻と観測地点が追加で必要になる理由を説明せよ。
\end{enumerate}

本章までに、データの読み込み、整形、計算、可視化、物理量の表現を扱った。同じ解析を別のデータへ安全に適用するには、処理を関数とモジュールへ分け、テスト可能な公開インタフェースを定める必要がある。次章では、そのためのパッケージ構成と依存関係の管理を扱う。
